\chapter{TypeScript + JSX prerequisites}

\section{The goal}
Relax.js supports writing UI in TSX/JSX with a React-like syntax.
But remember: JSX is only syntax.
It becomes function calls (or, in TypeScript's modern mode, calls to a small ``JSX runtime'' module).

In Relax, the important function is still \texttt{h()} from \texttt{src/h.ts}.
The TSX layer exists to make authoring more ergonomic and familiar.

This chapter answers four practical questions:

\begin{enumerate}
  \item What does TypeScript emit for TSX in this repo?
  \item Where is the TSX runtime implemented (and what does it do)?
  \item Which React-like conveniences are supported (and which are not)?
  \item How do tests lock the TSX contract in place?
\end{enumerate}

\section{How TSX works in this repository (TypeScript configuration)}
The TSX behavior in this repo is primarily defined in \texttt{tsconfig.json}:

\begin{itemize}
  \item \texttt{jsx: "react-jsx"}
  \item \texttt{jsxImportSource: "relax-jsx"}
\end{itemize}

This combination means:

\begin{itemize}
  \item You \emph{do not} configure a global JSX factory function.
  \item TypeScript emits imports from \texttt{relax-jsx/jsx-runtime} (and sometimes \texttt{relax-jsx/jsx-dev-runtime}).
\end{itemize}

So the contract becomes very clean:
``TSX syntax must compile into objects that Relax can mount/patch/destroy.''

\section{Two JSX styles youll see}

\subsection{Direct VDOM authoring}
You can write VNodes manually:

\begin{lstlisting}[style=relax,language=JavaScript]
import { h, hFragment } from 'relax.js'

const view = hFragment([
  h('h1', {}, ['Hello']),
  h('button', { on: { click: () => console.log('clicked') } }, ['Click'])
])
\end{lstlisting}

\subsection{React-like TSX authoring}
TypeScript can compile TSX into calls to a configured JSX runtime.
This repo uses a React-like authoring style where:

\begin{itemize}
  \item \texttt{className} maps to Relaxs \texttt{class}.
  \item \texttt{onClick} maps to Relaxs event prop form (internally becomes \texttt{on: \{ click: fn \}}).
  \item Fragments (\texttt{<>...</>}) are supported.
\end{itemize}

There's an important nuance here:

\begin{quote}
Relax supports both ``native'' Relax props (like \texttt{class} and \texttt{on: \{ click: ... \}}) \emph{and} React-like compatibility props (like \texttt{className} and \texttt{onClick}).
\end{quote}

This design keeps the underlying runtime simple (it only needs to understand one internal prop shape) while letting TSX authoring feel familiar.

\section{The Relax JSX runtime (the real implementation)}
The TSX runtime lives in \texttt{relax-jsx/jsx-runtime.ts}.
It's intentionally tiny.

At a high level it does four jobs:

\begin{enumerate}
  \item Normalizes \texttt{children} to an array.
  \item Applies React-like prop compatibility:
    \begin{itemize}
      \item \texttt{className -> class}
      \item \texttt{onClick/onInput/... -> on: \{ click/input/... \}}
    \end{itemize}
  \item Threads the optional \texttt{key} through into Relax props.
  \item Delegates to Relax's real VDOM constructors:\
    	exttt{h(type, props, children)} or \texttt{hFragment(children)}.
\end{enumerate}

\subsection{Fragments: why \texttt{Fragment} is exported}
In the modern TS transform, fragments compile to a special ``Fragment'' value.
Relax's JSX runtime exports:

\begin{itemize}
  \item \texttt{export const Fragment = hFragment as any}
\end{itemize}

Then \texttt{jsxImpl()} checks whether the element type is \texttt{Fragment} and returns \texttt{hFragment(children)}.
That makes fragment VNodes indistinguishable from ones created explicitly via \texttt{hFragment()}.

\subsection{Dev runtime: jsxDEV}
Tooling sometimes emits \texttt{jsxDEV} calls in development builds.
This repo includes \texttt{relax-jsx/jsx-dev-runtime.ts}, which simply re-exports \texttt{Fragment}, \texttt{jsx}, and \texttt{jsxs}, and provides \texttt{jsxDEV(...)} as an alias to \texttt{jsx(...)}.

This is not a performance hack; it's a compatibility bridge so Vite/Vitest can run TSX code in different modes.

\section{Tests as spec}
The definitive spec for supported TSX/JSX features lives in:

\begin{itemize}
  \item \texttt{src/__tests__/jsx-tsx.test.ts}
  \item \texttt{src/__tests__/jsx-react-syntax.test.tsx}
\end{itemize}

\subsection{Test 1: TSX can mount via the Relax JSX runtime}
	exttt{src/__tests__/jsx-tsx.test.ts} is the smallest end-to-end proof:

\begin{itemize}
  \item It imports and runs \texttt{examples/jsx/hello.tsx}.
  \item That example defines a component whose \texttt{render()} returns TSX.
  \item It mounts using \texttt{createApp()}, waits \texttt{nextTick()}, and returns \texttt{innerHTML}.
\end{itemize}

If this test fails, it usually means one of these broke:

\begin{itemize}
  \item TS didn't resolve \texttt{relax-jsx/jsx-runtime} the way we expect (config/alias problem).
  \item The TSX runtime produced VNodes in a shape \texttt{mountDOM()} doesn't handle.
\end{itemize}

\subsection{Test 2: React-like conveniences are supported}
	exttt{src/__tests__/jsx-react-syntax.test.tsx} is a feature matrix.
It's intentionally written as ``contracts'' in a comment at the top of the file.

It specifies (and proves) things like:

\begin{itemize}
  \item intrinsic elements accept \texttt{id}, \texttt{className}, and text children
  \item fragments (\texttt{<>...</>}) work
  \item \texttt{onClick} attaches a DOM click handler
  \item list rendering via \texttt{items.map(...)} works (including \texttt{key})
  \item patching TSX trees updates DOM without reconstructing the whole host
  \item TSX works for Relax components created via \texttt{defineComponent()}
\end{itemize}

This test file is a great example of how to design compatibility: pick the exact subset you want, and encode it as executable assertions.

\section{Whats not the same as React}
Even though the syntax is familiar, semantics are Relaxs:

\begin{itemize}
  \item Components are created via \texttt{defineComponent()}.
  \item The event binding model binds handlers to the host component instance.
  \item Slots in Relax are a component feature (external/default content), not React children.
\end{itemize}

\subsection{Practical limitations of the TSX subset}
This repo intentionally does \emph{not} try to be a full React clone.
Two examples:

\begin{itemize}
  \item The TSX runtime only rewrites ``event props'' when the value is a function.
    If you pass non-function values to \texttt{onClick}, you should expect it to behave as a normal prop.
  \item ``React children'' as a conceptual model is not the point.
    In Relax, children are just VNodes (or arrays of VNodes) passed to \texttt{h()}.
\end{itemize}

The tests are the source of truth for what is supported.
If you'd like to extend the supported subset, the safest workflow is:

\begin{enumerate}
  \item Add a failing test first in \texttt{src/__tests__/jsx-react-syntax.test.tsx}.
  \item Update \texttt{relax-jsx/jsx-runtime.ts} to normalize the new case.
  \item Ensure mount/patch/destroy tests still pass.
\end{enumerate}

\section{Mini-exercise}
Read \texttt{src/__tests__/jsx-react-syntax.test.tsx} and answer:

\begin{enumerate}
  \item How are fragments compiled?
  \item How are events passed?
  \item How do components appear in JSX and how do props flow?
\end{enumerate}

Then open \texttt{relax-jsx/jsx-runtime.ts} and locate the exact code that implements each answer.

\section{Online resources (optional, citation-ready)}

\begin{itemize}
  \item TypeScript handbook: JSX (explains \texttt{react-jsx} and \texttt{jsxImportSource}):
    \href{https://www.typescriptlang.org/docs/handbook/jsx.html}{TypeScript: JSX}.
  \item React JSX runtime reference (useful for understanding why \texttt{jsx}/\texttt{jsxs}/\texttt{jsxDEV} exist):
    \href{https://react.dev/reference/react/jsx-runtime}{React docs: jsx-runtime}.
  \item Vite + Vitest module resolution (useful when JSX runtime imports fail):
    \href{https://vite.dev/guide/features.html#jsx}{Vite: JSX},
    \href{https://vitest.dev/guide/}{Vitest Guide}.
  \item DOM events (what an \texttt{onClick} becomes):
    \href{https://developer.mozilla.org/en-US/docs/Web/API/Element/click_event}{MDN: click event}.
\end{itemize}
